\documentclass[12pt]{article}

\usepackage[margin=1.12in]{geometry}
\usepackage{amsmath,amssymb,amsthm,mathtools}
\usepackage{booktabs}
\usepackage{hyperref}

\newtheorem{theorem}{Theorem}
\newtheorem{proposition}{Proposition}
\theoremstyle{remark}
\newtheorem{remark}{Remark}

\newcommand{\R}{\mathbb R}
\newcommand{\Span}{\operatorname{span}}
\newcommand{\I}{\mathcal I}

\setlength{\emergencystretch}{3em}
\hypersetup{
  pdftitle={Probe Measurements Force Finite Centers Only Through Center-Separating Effects},
  pdfsubject={Fewster-Verch measurements, finite centers, local probes, tomography},
  pdfkeywords={Fewster-Verch, local measurement, finite centers, probe effects, tomography, AQFT}
}

\title{Probe Measurements Force Finite Centers\\
Only Through Center-Separating Effects}
\author{}
\date{Live review note of June 20, 2026}

\begin{document}
\maketitle

\begin{abstract}
Local probe-measurement frameworks in algebraic quantum field theory provide a
causal and covariant way to define induced system observables and instruments.
They are therefore a possible positive route out of finite-center
nonidentifiability.  This note isolates the finite-central boundary.  A
probe/instrument family reconstructs a finite central target exactly when the
induced system effects separate the corresponding central contrast fibers.
If all induced effects factor through a quotient, retraction, or
coarse-graining, then finite adaptive composition preserves that blindness.
Thus a measurement formalism alone does not force hidden finite central rank or
the predicate \(q=2\); the supplied probe effects must be informationally
complete for that target, or a physical quotient/exclusion axiom must remove
the hidden center.
\end{abstract}

\paragraph{Status.}
This is non-frozen external-review metadata for Team 1.  It is intended as a
small claim that a local-measurement, LCQFT, or AQFT reviewer can classify as
known, false, missing-hypothesis, or useful.  It does not modify any frozen
Team 1 manuscript, Lean archive, public mirror manifest, or SHA-256 package.

\section{Primary sources}

The boundary below was checked against the Fewster--Verch local measurement
framework \cite{FewsterVerchMeasurements}, the Fewster--Verch review of
measurement in QFT \cite{FewsterVerchReview}, the LCQFT/dynamical-locality
setting \cite{FewsterVerchDynamical}, and recent work on measurements in the
Fewster--Verch framework \cite{MandryschNavascues}.  These works are used here
as positive measurement constructions, not as claims that finite-center
separation is absent from physics.

\section{Finite central readout}

Let \(Z=\ell^\infty(\{1,\ldots,q\})\) be a finite center.  Its real contrast
space is
\[
  V_q=\left\{x=(x_1,\ldots,x_q)\in\R^q:\sum_i x_i=0\right\}.
\]
A declared family of probe preparations, couplings, and probe readouts induces
system effects.  Restricting those induced effects to \(Z\) gives vectors
\(e_\alpha\in\R^q\), hence central contrast functionals
\[
  x\longmapsto \langle e_\alpha,x\rangle .
\]
Let
\[
  N=\{x\in V_q:\langle e_\alpha,x\rangle=0
       \hbox{ for every supplied induced effect }\alpha\}.
\]

\begin{theorem}[probe effect-span criterion]
\label{thm:effect-span}
Let \(T:V_q\to X\) be a finite central target.  Then \(T\) is reconstructible
from the supplied probe-measurement data if and only if \(T\) is constant on
cosets of \(N\):
\[
  x-y\in N\quad\Longrightarrow\quad T(x)=T(y).
  \tag{1}
\]
In particular, the full finite central contrast is reconstructible exactly
when
\[
  \Span\{e_\alpha-\bar e_\alpha{\bf 1}\}_\alpha = V_q^*,
  \qquad
  \bar e_\alpha=\frac1q\sum_i(e_\alpha)_i.
  \tag{2}
\]
\end{theorem}

\begin{proof}
The supplied central measurement record is the linear map
\[
  R:V_q\to\R^\I,\qquad R(x)_\alpha=\langle e_\alpha,x\rangle,
\]
whose kernel is \(N\).  A target is a function of the measurement record
exactly when it is constant on the fibers of \(R\), which is (1).  Full
contrast reconstruction means that \(R\) is injective on \(V_q\), equivalently
that the restrictions of the induced effects span \(V_q^*\), which is (2).
\end{proof}

\section{Adaptive protocols}

\begin{proposition}[adaptivity does not create a missing central separator]
\label{prop:adaptive}
Suppose each stage of a finite adaptive protocol has induced central effects
that factor through a fixed central coarse-graining
\[
  C:\{1,\ldots,q\}\to\{1,\ldots,r\}.
\]
Then every joint outcome probability of the finite protocol factors through
the same coarse-graining.  Consequently no target that varies inside a
\(C\)-fiber is reconstructible from the protocol.
\end{proposition}

\begin{proof}
At each stage, conditioning on previous outcomes may choose a different
coupling, probe state, or probe effect, but by hypothesis the induced central
functional is constant on the fibers of \(C\).  Products, sums, and conditional
branches of functions constant on \(C\)-fibers remain constant on
\(C\)-fibers.  Induction over the finitely many stages shows that every joint
outcome probability is constant on \(C\)-fibers.  Any target that changes
within such a fiber therefore fails the fiber-constancy condition of
Theorem~\ref{thm:effect-span}.
\end{proof}

\section{Positive exits and referee question}

\begin{center}
\begin{tabular}{p{0.35\linewidth}p{0.50\linewidth}}
\toprule
Exit premise & Effect \\
\midrule
Center-separating induced effects & The probe observables span the central
contrast dual \(V_q^*\), so finite central states are tomographically visible.
\\
\addlinespace
Direct atomic central readout & Effects proportional to the minimal central
projections \(e_i\) separate every central atom. \\
\addlinespace
Factor, quotient, or inadmissibility axiom & The hidden finite center is
removed from the comparison class rather than reconstructed by measurement. \\
\addlinespace
Full physical tomography assumption & If the theorem assumes tomography for
the full physical algebra, finite central separation is part of the hypothesis.
\\
\bottomrule
\end{tabular}
\end{center}

The exact external question is:

\begin{quote}
Does the local probe-measurement setup in the target setting supply induced
system effects that separate the finite central contrasts, or does it only
provide a causal measurement framework whose chosen effects may still factor
through a quotient or retraction?
\end{quote}

If the first holds, the measurement scheme is a positive exit from the Team 1
finite-center obstruction.  If the second holds, the framework alone does not
force hidden finite central rank or \(q=2\).

\begin{thebibliography}{9}

\bibitem{FewsterVerchMeasurements}
C. J. Fewster and R. Verch,
\emph{Quantum fields and local measurements},
Commun. Math. Phys. 378 (2020), 851--889,
\href{https://arxiv.org/abs/1810.06512}{arXiv:1810.06512}.

\bibitem{FewsterVerchReview}
C. J. Fewster and R. Verch,
\emph{Measurement in Quantum Field Theory},
\href{https://arxiv.org/abs/2304.13356}{arXiv:2304.13356}.

\bibitem{FewsterVerchDynamical}
C. J. Fewster and R. Verch,
\emph{Dynamical locality and covariance: What makes a physical theory the same in all spacetimes?},
Ann. Henri Poincare 13 (2012), 1613--1674,
\href{https://arxiv.org/abs/1106.4785}{arXiv:1106.4785}.

\bibitem{MandryschNavascues}
J. Mandrysch and M. Navascu\'es,
\emph{Quantum Field Measurements in the Fewster--Verch Framework},
\href{https://arxiv.org/abs/2411.13605}{arXiv:2411.13605}.

\end{thebibliography}

\end{document}
